%document class
\documentclass[10pt,oneside]{book}
%%%% Page Info + Commands %%%%%{

%packages
\usepackage{pgfplots}
\pgfplotsset{compat=1.18}
\usepackage{amsmath}
\usepackage{amsfonts}
\usepackage{amsthm,letltxmacro,changepage}
\usepackage{enumerate}
\usepackage{enumitem}
\usepackage{changepage}

%This will shrink the page
\newcommand{\smallpage}{  \setlength \oddsidemargin{.5in}
            \setlength \textwidth{5in}
            \setlength \topmargin{0in}
            \setlength \textheight{9in}}


% Commands for sets of numbers
\newcommand{\Z}{\mathbb{Z}}\newcommand{\R}{\mathbb{R}}\newcommand{\N}{\mathbb{N}}

\newcommand{\n}{\vspace{0.15cm}\\}
\newcommand{\en}{\vspace{0.25cm}\\}

\newcommand{\gimage}[3]{
\begin{figure}[h]   
    \centering          
    \includegraphics[width=#2\textwidth]{#1}  
    \caption{#3}
    \label{fig:#1}
\end{figure}\\
}

\newcommand{\centerit}[1]{{\begin{center} #1 \end{center}}}
\newcommand{\ul}[1]{\underline{#1}}

%%%%%%%% command for graphics %%%%%%%%%%%%%
\usepackage{fancyhdr}
%}

%%%% Page Setup %%%%%%%
\smallpage\sloppy
\pagestyle{fancy}
\lhead{\large{\textbf{Day 2 - Proof? Proof?}}}
%\rfoot{\thepage/\pageref{LastPage} }
\setlength{\headheight}{10pt}


% Proof writing
\LetLtxMacro\oldproof\proof
\let\endoldproof\endproof
\renewenvironment{proof}[1][\proofname]
  {\vspace{0.2cm}\begin{adjustwidth}{2em}{2em}
   \oldproof[#1]}
  {\endoldproof
   \end{adjustwidth}}
\begin{document}

%
% Introductions
%
We just began the class talking about the free weekend, where people are going. Considerations include:
\begin{itemize}[itemsep=0pt]
    \item Vienna (\$30 - \$40) each way. Staying in the hostel costs \$30 per night.
    \item Stay in Budapest (no moneys)
    \item Whatever else you want to do
\end{itemize}
I think most people are going to Vienna and buying a two-day travel pass. 
\subsection*{Last Time}
Last class we covered:
\begin{itemize}[itemsep=0pt]
    \item Direct Arguments \& practice
    \item Cases (Sometime an argument needs several small cases)
    \item Division Algorithm, ``forms'' for integers. 
\end{itemize}
Everyone's jetlagged, and it's nighttime, and it can be frusterating when things don't click instantly. There were only 3 homework problems, but `they weren't the easiest three'.
\subsection*{Division Algorithm}
We primary covered the division algorithm last class, and how it is useful for solving various problems. The divsion algorithm takes the form:
\[n = qb + r\]
where $n$ is an integer. \n
Essentially, the division algoritm is saying, ``we can represent any integer in terms of this equation''. Although this sounds a little abstract, essentially imagine it like this:
\centerit{\fbox{Say I let $q = 3$. For any given integer $n$, I can represent it in the form $3b + r$.}}
Take the numbers 4, 8, and 9. I can represent them as:
\begin{align*}
    4 &= 3b + 1,\, \text{ for } b = 1\\
    8 &= 3b + 2,\, \text{ for } b = 2\\
    9 &= 3b + 0,\, \text{ for } b = 3
\end{align*}
Now, we don't just have to use $q = 3$, but is a convenient example, because if $q=3$, we have a maximum of 3 representations to cover a set of all integers: $3b + r$, for $r = \{0, 1, 2\}$. \n
For $q\ne 3$, things get more complicated, but it's still really managable. For say $q = 6$, we have six forms instead:
\centerit{
\begin{tabular}{c c c c c c}
    $6b + 0$ & $6b + 1$ & $6b + 2$ & $6b + 3$ & $6b+4$ & $6b+5$
\end{tabular}}
But you may be asking, \textit{how do we use this in a proof?} Well you kinda gotta figure this out by yourself, but you split your proof into \textbf{cases}. However, a hint is that you begin your proof like this:
\begin{proof}
    Let $n\in\Z$. By the \textbf{Division Algorithm}, $n$ has one of the forms $4b,\, 4b+1,\, 4b+2, $ or $4b+3$.\\
    \textbf{Case 1: } If $n= 4b$ ($b\in \Z$,) then $\ldots$
    \textbf{Case 2: } $\ldots$
\end{proof}
\pagebreak



\section*{New Material for Day 2}
After reviewing the division algorithm, we started govering the \textbf{Greatest Commond Divisor} (GCD).
\subsection*{Greatest Common Divisor}
Here's our working definition:
\begin{center}
    \fbox{\fbox{
        Given $a,b\in Z$, $\text{gcd}(a,b)$ is the largest integer $d$ for which $d\mid a$ and $d\mid b$. 
    }}
\end{center}
\footnotesize{*Note that, $d\mid a$ means that ``$d$ \text{divides} $a$'' }\normalsize\n
For example:
\centerit{
    \begin{tabular}{c c c}
        gcd(4,6) $=2$ & gcd$(4,5) = 1$ & gcd$(-8, -4) = 4$
    \end{tabular}
}
\subsection*{In-Class Work (pt. 1)}
\subsubsection*{Question 1}
\begin{enumerate}
    \item Find gcd$(24, 98)$
    It's 2
    \item Find the gcd$(1, 18)$
    It's 1
    \item Find the gcd$(-24, 98)$ It's 2
    \item Find the gcd$(1000, 1001)$ I'ts 1. 
\end{enumerate}
\subsubsection{Linear Combinations of Integers}
An interesting property of the greatest common divisor is that is is actually a \textbf{linear combination} of both numbers. Take gcd$(a,b)$, we can write:
\centerit{
    \fbox{gcd$(a,b) = x\cdot a + y\cdot b,\quad x,y\in\Z$. }
}
Little abstract and confusing, but imagine it like this. Let $a = 2, b = 3$. Their greatest common divisor is clearly 1. However, we can form the greatest common divisor from linearly combining $a$ and $b$:
\[-1 \cdot a + 1 \cdot b = 1\]
Wow! Look at that!
\subsubsection*{Question 2}
Find one pair of integers $x, y$ such that gcd$(a,b) = xa + yb$:
\begin{enumerate}[itemsep=0pt]
    \item $a = 10, b = 25:\quad \text{gcd}(10,25)=5,\quad x = -2, \;y = 1$
    \item $a = 10, b = 25:\quad \text{gcd}(24,98)=2,\quad x = -4, \;y = 1$
\end{enumerate}
\subsection*{Relative Prime}
You gotta know this definition. I don't know why, but you gotta know it:
\begin{center}
    \fbox{\fbox{Two integers $a, b$ are said to be \textit{relatively prime} if gcd$(a,b) = 1$}}
\end{center}

\pagebreak
\section*{Euclidean Aglorithm}
Now we need a way to consistenly solve for the gcd, and we'll use the euclidean algorithm to do so. Let $a$ and $b$ be the two integers we want to find the gcd for. We use the following code to find it:
\begin{verbatim}
while (b != 0) {
    int r = a % b;
    if (r == 0)
        break;
    a = b;
    b = r;
}
return b;
\end{verbatim} 
Now what is this algorithm actually doing? Well, let's break it down step by step. 
\begin{enumerate}[itemsep=0pt]
    \item First, find the remainder of $a \div b$. 
    \item \textbf{Case: }If our remainer is 0, the current value of $b$ is our gcd, and end here.
    \item Otherwise, set $a$ equal to the original $b$.
    \item Finally, set $b$ equal to the remainder.
    \item Repeat steps 1-4.
\end{enumerate}
We can see this in action with $a = 24$ and $b = 19$ (Ignore that they are relatively prime, this is for example purposes):
\begin{align*}
    24 / 19 &= 1(19) + 3\\
    19 / 3 &= 6(8) + 1\\
    3 / 1 &= 3(1) + 0
\end{align*}
Thus, because our last remainder is zero, we find that the remainder above that, 1, is our greatest common denominator, as desired. 

\section*{Diophantine Equations}
A Diophantine Equation is an equation that only concerns integer solutions to integer equations. For example, the equation:
\[ax+by =c\]
is a Diophantine equation if we restrict: $a,b,c,x,y\in \Z$. \n
A couple examples looks like:
\centerit{
    \begin{tabular}{c c c}
        $2x+4y =8$ & $3x + 5y = 13$ & $2 + 4y = 7$ 
    \end{tabular}
}
A interesting property of this particular subclass of equation, $ax + by = c$, is that:
\begin{center}
    \fbox{A solution to the diophantine eq. $ax + by = c$ exists if and only if gcd$(a,b) \mid c$}
\end{center}
It's that interesting? For example:
\[6x + 9y = 12\]
We know solutions exist because gcd$(6, 9) = 3$, and $3\mid 12$.\n
Let's go through an example problem.
\subsection*{Example Example Example Eggsample Example}
To solve $ax +by = c$:
\begin{itemize}
    \item Check if gcd$(a,b)\mid c$
    \item Solve $ax + by = \text{gcd}(a,b)$ using Euclidean Algorithm
    \item Multiply by appropriate multiplier to get a solution to $ax + by =c$ 
\end{itemize}
We didn't really finish this. 
\end{document}
